\section{Tren-ES}\label{tren-es}

\subsection{Supuesto elegido}\label{supuesto-elegido}

Hemos elegido realizar la práctica sobre un supuesto propio, distinto de
los propuestos. El objetivo consiste en construir un agregador que
integre información acerca de la situación ferroviaria nacional. El
sistema relacionará información sobre las estaciones, rutas y viajes en
tren por el territorio español. El objetivo final será realizar las
siguientes consultas:

\begin{enumerate}
\def\labelenumi{\arabic{enumi}.}
\tightlist
\item
  Estaciones de tren en poblaciones de menos de 10000 habitantes.
\item
  Viajes entre dos estaciones a menos de 30 km entre sí.
\item
  Poblaciones de entre 20000 y 100000 habitantes en las que haya menos
  de cinco viajes programados hoy.
\end{enumerate}

\subsection{Esquema mediador}\label{esquema-mediador}

Para modelar la información del sistema se ha definido un esquema
mediador compuesto por seis relaciones principales: \textbf{Estacion},
\textbf{Municipio}, \textbf{Viaje}, \textbf{Parada}, \textbf{Ruta} y
\textbf{Distancia}. Este esquema permite representar tanto la estructura
geográfica del sistema ferroviario como la planificación de trayectos y
viajes.

Las relaciones del esquema mediador son las siguientes:

\begin{itemize}
\tightlist
\item
  Estacion(\textbf{id}, nombre, \emph{municipio}, latitud, longitud)
\item
  Municipio(nombre, n\_habitantes, \textbf{id}, latitud, longitud,
  provincia, CCAA)
\item
  Viaje(\textbf{id}, \emph{ruta}, fecha, horario)
\item
  Parada(\textbf{\emph{ruta}}, \textbf{\emph{estacion}}, n\_secuencia,
  km\_origen)
\item
  Ruta(\textbf{id}, \emph{origen}, \emph{destino}, tipo)
\item
  Distancia(\textbf{\emph{estacion1}}, \textbf{\emph{estacion2}},
  distancia)
\end{itemize}

A continuación se describe brevemente cada una de las tablas y los
atributos que la componen.

\subsubsection{Tabla Estacion}\label{tabla-estacion}

La tabla \textbf{Estacion} guarda la información de las estaciones de
tren registradas en el sistema. Cada estación está asociada a un
municipio concreto.

Sus atributos son:

\begin{itemize}
\tightlist
\item
  \textbf{id}: clave primaria de la tabla. Identifica de forma única
  cada estación.
\item
  \textbf{nombre}: nombre de la estación.
\item
  \textbf{municipio}: clave foránea que referencia al atributo
  \textbf{id} de la tabla \textbf{Municipio}. Indica el municipio en el
  que se encuentra situada la estación.
\item
  \textbf{latitud}: coordenada geográfica de latitud de la estación.
\item
  \textbf{longitud}: coordenada geográfica de longitud de la estación.
\end{itemize}

\subsubsection{Tabla Municipio}\label{tabla-municipio}

La tabla \textbf{Municipio} almacena la información básica de cada
municipio incluido en el sistema. En ella se recogen tanto datos
identificativos como datos demográficos y geográficos.

Sus atributos son:

\begin{itemize}
\tightlist
\item
  \textbf{id}: clave primaria de la tabla. Identifica de manera única
  cada municipio.
\item
  \textbf{nombre}: nombre del municipio.
\item
  \textbf{n\_habitantes}: número de habitantes del municipio.
\item
  \textbf{latitud}: coordenada geográfica de latitud del municipio.
\item
  \textbf{longitud}: coordenada geográfica de longitud del municipio.
\item
  \textbf{provincia}: provincia a la que pertenece el municipio.
\item
  \textbf{CCAA}: comunidad autónoma a la que pertenece el municipio.
\end{itemize}

\subsubsection{Tabla Viaje}\label{tabla-viaje}

La tabla \textbf{Viaje} almacena los viajes concretos programados en el
sistema. Mientras que una ruta define un recorrido general, un viaje
representa una realización específica de dicha ruta en una fecha y hora
determinadas.

Sus atributos son:

\begin{itemize}
\tightlist
\item
  \textbf{id}: clave primaria de la tabla. Identifica de forma única
  cada viaje.
\item
  \textbf{ruta}: clave foránea que referencia al atributo \textbf{id} de
  la tabla \textbf{Ruta}. Indica qué ruta realiza el viaje.
\item
  \textbf{fecha}: fecha en la que está programado el viaje.
\item
  \textbf{horario}: hora programada para la salida del viaje.
\end{itemize}

\subsubsection{Tabla Parada}\label{tabla-parada}

La tabla \textbf{Parada} almacena las estaciones por las que pasa una
ruta, guardando el orden de llegada planificado en la ruta a la que
pertenece. Esta tabla NO tiene en cuenta lla existencia de rutas
circulares (que por otra parte creemos que no existen en recorridos de
alta y media distancia). Consideramos que las rutas tienen un origen y
un destino, con una serie de paradas intermedias.

Sus atributos son:

\begin{itemize}
\tightlist
\item
  \textbf{ruta}: clave foránea que referencia al atributo \textbf{id} de
  la tabla \textbf{Ruta}. Forma parte de la clave primaria compuesta de
  la tabla.
\item
  \textbf{estacion}: clave foránea que referencia al atributo
  \textbf{id} de la tabla \textbf{Estacion}. Forma parte de la clave
  primaria compuesta de la tabla.
\item
  \textbf{n\_secuencia}: indica el orden de la parada dentro de la ruta.
\item
  \textbf{km\_origen}: distancia en kilómetros desde el origen de la
  ruta hasta esa parada.
\end{itemize}

La clave primaria de \textbf{Parada} es compuesta, formada por los
atributos \textbf{ruta} y \textbf{estacion}.

\subsubsection{Tabla Ruta}\label{tabla-ruta}

La tabla \textbf{Ruta} representa los recorridos ferroviarios entre
estaciones. Cada ruta conecta una estación de origen con una estación de
destino y permite clasificar el trayecto según su tipo.

Sus atributos son:

\begin{itemize}
\tightlist
\item
  \textbf{id}: clave primaria de la tabla. Identifica de manera única
  cada ruta.
\item
  \textbf{origen}: clave foránea que referencia al atributo \textbf{id}
  de la tabla \textbf{Estacion}. Indica la estación de origen de la
  ruta.
\item
  \textbf{destino}: clave foránea que referencia al atributo \textbf{id}
  de la tabla \textbf{Estacion}. Indica la estación de destino de la
  ruta.
\item
  \textbf{tipo}: tipo de ruta o servicio ferroviario asociado.
\end{itemize}

\subsubsection{Tabla Distancia}\label{tabla-distancia}

La tabla \textbf{Distancia} almacena la distancia entre pares de
estaciones. Esta relación se utiliza para facilitar consultas sobre
proximidad geográfica entre estaciones.

Sus atributos son:

\begin{itemize}
\tightlist
\item
  \textbf{estacion1}: clave foránea que referencia al atributo
  \textbf{id} de la tabla \textbf{Estacion}. Forma parte de la clave
  primaria compuesta.
\item
  \textbf{estacion2}: clave foránea que referencia al atributo
  \textbf{id} de la tabla \textbf{Estacion}. Forma parte de la clave
  primaria compuesta.
\item
  \textbf{distancia}: distancia entre ambas estaciones.
\end{itemize}

La clave primaria de \textbf{Distancia} es compuesta, formada por
\textbf{estacion1} y \textbf{estacion2}.

\subsection{Diagrama entidad-relación del
mediador}\label{diagrama-entidad-relaciuxf3n-del-mediador}

\begin{figure}
\centering
\includegraphics{./img/mediador.jpeg}
\caption{Esquema del mediador}
\end{figure}

\subsection{Formulación SQL de las
consultas}\label{formulaciuxf3n-sql-de-las-consultas}

Podemos formular las consultas de la siguiente manera

\begin{enumerate}
\def\labelenumi{\arabic{enumi}.}
\tightlist
\item
  \textbf{Estaciones de tren en poblaciones de menos de 10000
  habitantes.}
\end{enumerate}

\begin{Shaded}
\begin{Highlighting}[]
\KeywordTok{SELECT}\NormalTok{ E.}\KeywordTok{id}\NormalTok{, E.nombre }\KeywordTok{AS}\NormalTok{ estacion, M.nombre }\KeywordTok{AS}\NormalTok{ municipio}
    \KeywordTok{FROM}\NormalTok{ Estacion }\KeywordTok{AS}\NormalTok{ E}
    \KeywordTok{INNER} \KeywordTok{JOIN}\NormalTok{ Municipio }\KeywordTok{AS}\NormalTok{ M}
        \KeywordTok{ON}\NormalTok{ E.municipio }\OperatorTok{=}\NormalTok{ M.}\KeywordTok{id}
    \KeywordTok{WHERE}\NormalTok{ M.n\_habitantes }\OperatorTok{\textless{}} \DecValTok{10000}\NormalTok{;}
\end{Highlighting}
\end{Shaded}

\begin{enumerate}
\def\labelenumi{\arabic{enumi}.}
\setcounter{enumi}{1}
\tightlist
\item
  \textbf{Viajes entre dos estaciones a menos de 30 km entre sí.}
\end{enumerate}

Para la solución a esta consulta es importante ver a qué nos referimos
con \emph{distancia}: es la distancia euclídea directa entre las
estaciones origen y destino a partir de su longitud y latitud. No tiene
en cuenta el recorrido que haga, pudiendo ser este mayor que 30 km.

\begin{Shaded}
\begin{Highlighting}[]
\KeywordTok{SELECT}\NormalTok{ V.}\KeywordTok{id}
    \KeywordTok{FROM}\NormalTok{ Viaje }\KeywordTok{AS}\NormalTok{ V}
    \KeywordTok{INNER} \KeywordTok{JOIN}\NormalTok{ Ruta }\KeywordTok{AS}\NormalTok{ R}
        \KeywordTok{ON}\NormalTok{ V.ruta }\OperatorTok{=}\NormalTok{ R.}\KeywordTok{id}
    \KeywordTok{INNER} \KeywordTok{JOIN}\NormalTok{ Distancia }\KeywordTok{AS}\NormalTok{ D}
        \KeywordTok{ON}\NormalTok{ (R.origen }\OperatorTok{=}\NormalTok{ D.estacion1 }\KeywordTok{AND}\NormalTok{ R.destino }\OperatorTok{=}\NormalTok{ D.estacion2) }\KeywordTok{OR}
\NormalTok{        (R.origen }\OperatorTok{=}\NormalTok{ D.estacion2 }\KeywordTok{AND}\NormalTok{ R.destino }\OperatorTok{=}\NormalTok{ D.estacion1)}
    \KeywordTok{WHERE}\NormalTok{ D.distancia }\OperatorTok{\textless{}} \DecValTok{30}\NormalTok{;}
\end{Highlighting}
\end{Shaded}

\begin{enumerate}
\def\labelenumi{\arabic{enumi}.}
\setcounter{enumi}{2}
\tightlist
\item
  \textbf{Poblaciones de entre 20000 y 100000 habitantes en las que haya
  entre uno y cinco viajes programados hoy.}
\end{enumerate}

Esta consulta en principio iba a ser **Poblaciones \ldots{} haya menos
de cinco viajes programados hoy. El problema con esta consulta así
formulada es que es necesario hacer un \emph{OUTER JOIN} para obtener
aquellas poblaciones en las que haya 0 viajes programados, lo que daba
problemas al poner las consultas en forma conjuntiva, por lo que hemos
decidido realizar esta simplificación.

\begin{Shaded}
\begin{Highlighting}[]
\KeywordTok{SELECT}
\NormalTok{    M.}\KeywordTok{id}\NormalTok{,}
\NormalTok{    M.nombre,}
    \FunctionTok{COUNT}\NormalTok{(}\KeywordTok{DISTINCT}\NormalTok{ V.}\KeywordTok{id}\NormalTok{) }\KeywordTok{AS}\NormalTok{ num\_viajes\_hoy}
\KeywordTok{FROM}
\NormalTok{    Municipio }\KeywordTok{AS}\NormalTok{ M}
    \KeywordTok{INNER} \KeywordTok{JOIN}\NormalTok{ Estacion }\KeywordTok{AS}\NormalTok{ E }\KeywordTok{ON}\NormalTok{ E.municipio\_id }\OperatorTok{=}\NormalTok{ M.}\KeywordTok{id}
    \KeywordTok{INNER} \KeywordTok{JOIN}\NormalTok{ Parada }\KeywordTok{AS}\NormalTok{ P }\KeywordTok{ON}\NormalTok{ P.estacion }\OperatorTok{=}\NormalTok{ E.}\KeywordTok{id}
    \KeywordTok{INNER} \KeywordTok{JOIN}\NormalTok{ Viaje }\KeywordTok{AS}\NormalTok{ V }\KeywordTok{ON}\NormalTok{ V.ruta }\OperatorTok{=}\NormalTok{ P.ruta}
\KeywordTok{WHERE}
\NormalTok{    M.n\_habitantes }\KeywordTok{BETWEEN} \DecValTok{20000} \KeywordTok{AND} \DecValTok{100000}
    \KeywordTok{AND}\NormalTok{ V.fecha }\OperatorTok{=} \FunctionTok{CURRENT\_DATE}
\KeywordTok{GROUP} \KeywordTok{BY}
\NormalTok{    M.}\KeywordTok{id}\NormalTok{,}
\NormalTok{    M.nombre}
\KeywordTok{HAVING}
    \FunctionTok{COUNT}\NormalTok{(}\KeywordTok{DISTINCT}\NormalTok{ V.}\KeywordTok{id}\NormalTok{) }\KeywordTok{BETWEEN} \DecValTok{1} \KeywordTok{AND} \DecValTok{3}\NormalTok{;}
\end{Highlighting}
\end{Shaded}

\subsection{Captura de una instancia}\label{captura-de-una-instancia}

En este apartado se muestran las capturas de instancias específicas del
esquema mediador:

\includegraphics{./img/instanciaRuta.jpeg}
\includegraphics{./img/instanciaDistancia.jpeg}
\includegraphics{./img/instanciaEstacion.jpeg}
\includegraphics{./img/instanciaParada.jpeg}
\includegraphics{./img/instanciaViaje.jpeg}

\section{Esquemas Origen}\label{esquemas-origen}

En este documento se muestran los esquemas de las fuentes de datos
origen de nuestro sistema integrador Tren-ES. Las fuentes de datos son
las siguientes:

\begin{itemize}
\tightlist
\item
  Wikidata. Para obtener información de los municipios, como su nombre,
  código del INE, coordenadas geográficas o la población. Para extraer
  estos datos utilizaremos un \emph{punto sparql}
\item
  INE. Para conectar los códigos de los municipios con las provincias y
  comunidades autónomas. Para extraer estos datos utilizaremos la
  \emph{API oficial}.
\item
  Renfe. Dos distintas:

  \begin{itemize}
  \tightlist
  \item
    Renfe Data. Obtenemos información a partir de la \emph{API de Renfe}
    de todas las estaciones que existen en España.
  \item
    Renfe horarios. Un pequeño formulario desde el que se puede obtener
    la información de todos los horarios de tramos (Estación Origen,
    Estación Destino) disponibles en la web de renfe y las rutas
    correspondientes. Para extraer estos datus utilizaremos \emph{web
    scraping}.
  \end{itemize}
\item
  ADIF. Información en tiempo real de las salidas y llegadas en cada
  estación. Nos permite contrastar los horarios planificados con los
  reales. Para extraer estos datus utilizaremos \emph{web scraping}.
\end{itemize}

\subsection{Forma conjuntiva de las tablas
origen}\label{forma-conjuntiva-de-las-tablas-origen}

\begin{itemize}
\tightlist
\item
  wikidata\_MUNICIPIO(codigo INTEGER, label VARCHAR(255), poblacion
  INTEGER, coordenadas POINT) Tabla de datos extraidos de Wikidata. La
  tabla está formada por un campo \texttt{codigo} que contiene el
  identificador INE del municipio, un campo \texttt{label} con el nombre
  del municipio en texto legible, la columna \texttt{poblacion} que
  almacena el número de habitantes de ese municipio y las
  \texttt{coordenadas}, una tupla de tipo Point que almacena las
  coordenadas geográficas de dicho municipio.
\item
  INE\_MUNICIPIO(cmun INTEGER, cpro INTEGER, dc INTEGER, nombre
  VARCHAR(100)) Datos oficiales sobre municipios ofrecidos por el INE.
  Formados por un código de municipio \texttt{cmun}, que identifica a un
  municipio \emph{dentro} de una provincia, el código de la provincia a
  la que pertenece, \texttt{cpro}, un dígito de control (que no nos
  genera interés) \texttt{dc} y el nombre en texto legible,
  \texttt{nombre}.
\item
  INE\_PROVINCIA(cpro INTEGER, nombre VARCHAR(100), codauto INTEGER,
  ccaa VARCHAR(100)) Datos oficiales sobre provincias ofrecidos por el
  INE. Es una tabla que cruza provincias con su respectiva comunidad
  autónoma. Formada por la información de provincia, \texttt{cpro} y
  \texttt{nombre}, código de la provincia y nombre en texto legible
  respectivamente, e información sobre la Comunidad Autónoma a la que
  pertenece, \texttt{codauto} y \texttt{ccaa}, código de la Comunidad
  Autónoma y nombre en texto legible respectivamente.
\item
  data\_renfe\_ESTACION(id INTEGER, codigo INTEGER, descripcion
  VARCHAR(255), latitud NUMERIC(9,6), longitud NUMERIC(9,6), direccion
  VARCHAR(255), c\_p VARCHAR(5), poblacion VARCHAR(100), provincia
  VARCHAR(100), pais VARCHAR(100)) Este esquema, se ofrece en la sección
  de datos abiertos de Renfe. Contiene información sobre todas las
  estaciones de trenes que hay en España. Se compone por, un \texttt{id}
  de asignación secuencial, un código de estación \texttt{codigo}, el
  nombre en texto legible guardado en \texttt{descripción}, las
  coordenadas geográficas por separado en \texttt{latitud} y
  \texttt{longitud}, una dirección guardada en la columna
  \texttt{direccion} y el código postal asociado \texttt{c\_p} y
  finalmente datos territoriales de la ubicación, \texttt{población},
  \texttt{provincia} y \texttt{pais}.
\item
  horarios\_renfe\_HORARIO(recorrido VARCHAR(255), salida TIME, llegada
  TIME, duracion INTERVAL) Datos extraídos de un formulario de la página
  de Renfe. El horario son viajes que van desde una estación de origen
  hasta una de destino. \texttt{recorrido} es la ruta que realiza ese
  viaje, \texttt{salida} y \texttt{llegada} son las estaciones que une
  dicho viaje (y que son paradas de la ruta). Duración es el tiempo en
  horas y minutos que tarda en completarse el viaje.
\item
  horarios\_renfe\_RUTA(codigo INTEGER, tipo VARCHAR(20), paradas TEXT)
  Se extraen del mismo formulario, contiene información sobre una ruta
  en específico. Las rutas son conjuntos de estaciones que se recorren
  en un orden específico. Una ruta tiene un \texttt{codigo} que la
  identifica frente al resto, y un \texttt{tipo} de tren que la recorre
  (MD, AV, AVLO, \ldots) y \texttt{paradas} representa la lista de
  paradas. Sabemos que no es una representación normalizada, pero nos
  parecía que añadía mucha complejidad a representar los orígenes, y por
  motivos de simpleza lo hemos hecho así. En el esquema mediador se
  trata y soluciona este problema.
\item
  adif\_SALIDAS(hora TIME, destino VARCHAR(255), tren VARCHAR(100), via
  INTEGER, estacion VARCHAR(255)) Información accesible desde las
  páginas de datos en tiempo real sobre las estaciones españolas.
  Contiene información sobre las salidas en tiempo real. \texttt{hora}
  es el cuándo sale el tren de la estación, el campo \texttt{destino}
  contiene la estación destino de dicha salida, \texttt{tren} almacena
  tanto la ruta de dicho viaje que sale como el tipo de tren que la
  atiende, una columna \texttt{via} desde dónde es dicha salida y
  \texttt{estacion} la estación de la que provienen los datos.\\
\item
  adif\_LLEGADAS(hora TIME, origen VARCHAR(255), tren VARCHAR(100), via
  INTEGER, estacion VARCHAR(255)) La entidad de llegadas contiene la
  misma información que salidas, pero intercambiando cualquier dato por
  los correspondientes a las llegadas.
\end{itemize}

\includegraphics{./img/origen1.jpeg}
\includegraphics{./img/origen2.jpeg}
\includegraphics{./img/origen3.jpeg}
\includegraphics{./img/origen4.jpeg}

A partir del mediador y los esquemas de las fuentes de datos podemos
acabar con los siguientes esquemas GAV:

\subsubsection{Municipio}\label{municipio}

Es necesario cruzar los datos de los datos de \emph{wikidata} y los del
\emph{INE} para poder obtener de forma correcta los códigos de provincia
y comunidad autónoma.

\texttt{Municipio(id,\ nombre,\ n\_habitantes,\ latitud\_m,\ longitud\_m,\ provincia,\ CCAA)}
\(\supseteq\)

\begin{verbatim}
wikidata_MUNICIPIO(id, nombre ,n_habitantes, coordenadas), latitud_m=latitud(coordenadas), 
longitud_m=longitud(coordenadas), INE_MUNICIPIO(cmun, cpro, v1, v2), id==cmun+cpro, 
INE_PROVINCIA(cpro, provincia, v3, CCAA)
\end{verbatim}

\subsubsection{Estacion}\label{estacion}

Hay que tener en cuenta que solo debemos obtener estaciones de España.

\texttt{Estacion(id,\ municipio,\ nombre,\ latitud\_e,\ longitud\_e)}
\(\supseteq\)

\begin{verbatim}
data_renfe_ESTACION(v1, id, nombre, latitud_e, longitud_e, v2, v3, nombre_mun, v4, v5), 
wikidata_MUNICIPIO(municipio, nombre_mun ,v6, pais), pais==España
\end{verbatim}

\subsubsection{Distancia}\label{distancia}

En este caso solo guardamos las distancias entre las estaciones de forma
unidireccional, ya que la distancia entre una estación y otra es la
misma que entre la segunda y la primera, y así evitamos información
repetida.

\texttt{Distancia(estacion1,\ estacion2,\ distancia)} \(\supseteq\)

\begin{verbatim}
data_renfe_ESTACION(v1, estacion1, v2, latitud_1, longitud_1, v3, v4, v5, v6, pais), 
data_renfe_ESTACION(v8, estacion2, v9, latitud_2, longitud_2, v10, v11, v12, v13, pais), 
distancia=dist(latitud_1, latitud_2, longitud_1, longitud_2), estacion1 < estacion2, 
pais==España
\end{verbatim}

\subsubsection{Ruta y Parada}\label{ruta-y-parada}

En las fuentes de datos obtenemos desde \emph{horarios\_renfe\_RUTA} una
lista de paradas. Para poder extraer cada parada por separado debemos
normalizar esta lista, por lo que nos ayudamos de una vista auxiliar.
Esta vista no se encuentra en los archivos SQL en este hito.

\paragraph{Vista auxiliar: Parada\_aux}\label{vista-auxiliar-parada_aux}

Esta vista auxiliar no se encuentra en el esquema mediador, utilizaremos
una función para normalizar la lista y poder extraer las paradas:

\texttt{Parada\_aux(codigo\_ruta,\ num\_secuencia,\ nombre\_parada)}\(\supseteq\)

\begin{verbatim}
horarios_renfe_RUTA(lista_paradas, codigo_ruta), 
(num_secuencia, nombre_parada) =unlist(lista_paradas)
\end{verbatim}

A partir de esta vista auxiliar podemos definir correctamente parada:

\texttt{Parada(ruta,\ estacion,\ num\_secuencia,\ km\_origen)}
\(\supseteq\)

\begin{verbatim}
Parada_aux(ruta, num_secuencia, nombre_parada), 
data_renfe_ESTACION(v1, estacion, nombre_parada, v3, v4, v5, v6, v7, v8, pais), pais==España
\end{verbatim}

Para definir la ruta volvemos a necesitar funciones externas. Una ruta
se define por un identificador propio y una tupla origen y destino de
estaciones de la ruta. Para esto debemos obtener los números de
secuencia mínimo y máximo de cada entrada de \emph{Parada\_aux}.

\texttt{Ruta(id,\ origen,\ destino,\ tipo)} \(\supseteq\)

\begin{verbatim}
Parada_aux(id, num_secuencia, nombre_parada), nombre_origen=getMinSecuencia(num_secuencia), 
nombre_destino=getMaxSecuencia(num_secuencia), data_renfe_ESTACION(v1, origen, nombre_origen, 
v3, v4, v5, v6, v7, v8, pais), data_renfe_ESTACION(v10, destino, nombre_destino, v11, v12, 
v13, v14, v15, v16, pais), pais==España
\end{verbatim}

\subsubsection{Viaje}\label{viaje}

El horario se refiere a la hora de llegada en la estación destino del
viaje específico, recogido de la web de \emph{ADIF}.

\texttt{Viaje(id,\ ruta,\ fecha,\ horario)} \(\supseteq\)

\begin{verbatim}
Parada_aux(ruta, num_secuencia, nombre_parada), nombre_origen=getMinSecuencia(num_secuencia), 
data_renfe_ESTACION(v1, estacion, nombre_origen, v3, v4, v5, v6, v7, v8, pais), 
adif_SALIDAS(fecha, horario, v10, v11, estacion), id=secuencia(), pais==España
\end{verbatim}

No se han considerado simplificaciones para ninguna tabla. Para la tabla
Municipio no se puede simplificar ya que es necesario hacer el cruce de
tablas con el INE para obtener correctamente el código de provincia y
Comunidad Autónoma. Para el resto de tablas no vemos simplificación
posible.

A partir del esquema mediador y los esquemas de las fuentes de datos,
podemos obtener los siguientes esquemas LAV:

\subsubsection{wikidata\_MUNICIPIO}\label{wikidata_municipio}

Necesitamos concatenar las coordenadas de latitud y longitud extraídas
en los esquemas GAV

\texttt{wikidata\_MUNICIPIO(codigo,\ label,\ poblacion,\ coordenadas)}\(\subseteq\)

\begin{verbatim}
Municipio(codigo, label, poblacion, latitud, longitud, v1, v2), 
coordenadas = coords(latitud, longitud)
\end{verbatim}

\subsubsection{INE\_MUNICIPIO}\label{ine_municipio}

Como no extraemos el dígito de control dc, no podemos obtenerlo de los
esquemas del mediador.

\texttt{INE\_MUNICIPIO(cmun,\ cpro,\ dc,\ nombre)}\(\subseteq\)

\begin{verbatim}
Municipio(id, nombre, v1, v2, v3, v4, v5), id == cmun + cpro
\end{verbatim}

\subsubsection{INE\_PROVINCIA}\label{ine_provincia}

Como no extraemos el código de CCAA, no podemos obtenerlo de los
esquemas del mediador.

\texttt{INE\_PROVINCIA(cpro,\ nombre,\ codauto,\ ccaa)}\(\subseteq\)

\begin{verbatim}
Municipio(id, nombre, v1, v2, v3, v4, v5), id == cmun + cpro
\end{verbatim}

\subsubsection{data\_renfe\_ESTACION}\label{data_renfe_estacion}

Como estación venía de un formulario con varios campos que no extraemos,
no podemos reconstruir exactamente la fuente completa a partir de
nuestros esquemas.

\texttt{data\_renfe\_ESTACION(v1,\ id,\ descripcion,\ latitud,\ longitud,\ v2,\ c\_p,\ poblacion,\ provincia,\ pais)}\(\subseteq\)

\begin{verbatim}
Estacion(id, id_mun, descripcion, latitud, longitud), 
Municipio(id_mun, poblacion, v3, v4, v5, provincia, v6)
\end{verbatim}

\subsubsection{horarios\_renfe\_RUTA}\label{horarios_renfe_ruta}

Al contrario que en los esquemas GAV, debemos crear una lista a partir
de nuestros esquemas.

\texttt{horarios\_renfe\_RUTA(codigo,\ paradas)}\(\subseteq\)

\begin{verbatim}
Parada(codigo, estacion, n_secuencia, v4), paradas = list(estacion, n_secuencia)
\end{verbatim}

\subsubsection{adif\_SALIDAS}\label{adif_salidas}

No podemos conseguir los campos tren y via desde nuestros esquemas, por
lo que la tabla no se puede construir por completo.

\texttt{adif\_SALIDAS(hora,\ destino,\ tren,\ via,\ estacion)}\(\subseteq\)

\begin{verbatim}
Viaje(id_viaje, ruta, fecha, hora), Ruta(ruta, origen, id_estacion_destino, tipo), 
Estacion(id_estacion_destino, estacion, destino, v2, v3)
\end{verbatim}

Al igual que en los esquemas GAV, no se ha considerado ninguna
simplificación.

\section{Reformulación GAV y LAV de las
Consultas}\label{reformulaciuxf3n-gav-y-lav-de-las-consultas}

En esta sección vamos a hacer la reformulación GAV y LAV de las
consultas. Para la reformulación GAV emplearemos la descomposición de
consultas y para la reformulación LAV utilizaremos el algoritmo de los
\emph{Buckets}.

\subsection{Consulta 1}\label{consulta-1}

\begin{itemize}
\tightlist
\item
  Estaciones de tren en poblaciones de menos de 10000 habitantes.
\end{itemize}

La consulta se puede escribir sobre el esquema mediador como:

\begin{verbatim}
Q1(id, nombre_estacion, nombre_municipio):- Estacion(id, id_mun, nombre_estacion, 
lat_e, lon_e), Municipio(nombre_municipio, habitantes, id_mun, lat_m, lon_m, prov, ccaa), 
hab<10000
\end{verbatim}

\subsubsection{Reformulación GAV}\label{reformulaciuxf3n-gav}

Para hacer la reformulación GAV basta con sustituir los átomos en la
consulta global por sus respectivas definiciones en función de las
fuentes, a partir de la formulación GAV de las fuentes vistas en
apartados anteriores.

Para el átomo \emph{Estacion}, sustituimos por su correspondiente
formulación GAV:
\texttt{Estacion(id,\ municipio,\ nombre,\ latitud\_e,\ longitud\_e)}
\(\subseteq\)

\begin{verbatim}
data_renfe_ESTACION(v1, id, nombre, latitud_e, longitud_e, v2, v3, nombre_mun, v4, v5), 
wikidata_MUNICIPIO(municipio, nombre_mun ,v6, pais), pais==España
\end{verbatim}

E igualmente con \emph{Municipio}, que queda:
\texttt{Municipio(id,\ nombre,\ n\_habitantes,\ latitud\_m,\ longitud\_m,\ provincia,\ CCAA)}
\(\subseteq\)

\begin{verbatim}
wikidata_MUNICIPIO(id, nombre ,n_habitantes, coordenadas), latitud_m=latitud(coordenadas), 
longitud_m=longitud(coordenadas), INE_MUNICIPIO(cmun, cpro, v1, v2), id==cmun+cpro, 
INE_PROVINCIA(cpro, provincia, v3, CCAA)
\end{verbatim}

Incluimos dichas formulaciones en la consulta global, añadiendo la
restricción de habitantes, que da como resultado:

\begin{verbatim}
Q1'(id, nombre_estacion, nombre_municipio):- data_renfe_ESTACION(v1, id, nombre_estacion, 
lat_e, lon_e, v2, v3, nom_mun, v4, v5), wikidata_MUNICIPIO(id_mun, nombre_municipio, v6, v7), 
wikidata_MUNICIPIO(id_mun, nombre_municipio, hab, coord), lat_m=latitud(coord), 
lon_m=longitud(coord), INE_MUNICIPIO(cmun, cpro, v8, v9), 
INE_PROVINCIA(cpro, prov, v10, ccaa), 
id_mun=cmun+cpro, hab<10000
\end{verbatim}

Por último, podemos ver que el átomo \emph{wikidata\_MUNICIPIO} se
encuentra repetido, por lo que podríamos pensar en simplificar para
eliminar redundancias. Podríamos hacer mappings de contención para ver
equivalencias, pero es evidente que, dado que v6 y v7 no se utilizan,
podemos eliminar el primer átomo, quedando la consulta final como:

\begin{verbatim}
Q1'(id, nombre_estacion, nombre_municipio):- data_renfe_ESTACION(v1, id, nombre_estacion, 
lat_e, lon_e, v2, v3, nom_mun, v4, pais), wikidata_MUNICIPIO(id_mun, nombre_municipio, 
hab, coord), lat_m=latitud(coord), lon_m=longitud(coord), 
INE_MUNICIPIO(cmun, cpro, v8, v9), INE_PROVINCIA(cpro, prov, v10, ccaa), 
id_mun=cmun+cpro, hab<10000, pais=España
\end{verbatim}

\subsubsection{Reformulación LAV}\label{reformulaciuxf3n-lav}

Para hacer la reformulación LAV debemos utilizar el algoritmo basado en
buckets.

\begin{enumerate}
\def\labelenumi{\arabic{enumi}.}
\tightlist
\item
  Creamos y llenamos los buckets: como hay dos átomos en la consulta,
  creamos dos buckets. Para que una vista forme parte de un bucket es
  necesario que cumpla tres condiciones:

  \begin{enumerate}
  \def\labelenumii{\alph{enumii}.}
  \tightlist
  \item
    Un átomo de la vista y el átomo g para el cual construimos el bucket
    afectan a la misma relación.
  \item
    Los predicados (condiciones) de la vista y los de la consulta Q son
    mutuamente satisfacibles.
  \item
    Si en g aparece alguna variable cabecera de la consulta Q, entonces
    también debe aparecer en la cabecera de la vista.
  \end{enumerate}
\end{enumerate}

Siguiendo las normas de construcción de buckets, para el primer átomo,
su bucket estará compuesto únicamente por \emph{wikidata\_ESTACION}, que
se añade al bucket. El segundo bucket estará compuesto únicamente por la
fuente \emph{wikidata\_MUNICIPIO}, que se añade al segundo bucket.

\begin{enumerate}
\def\labelenumi{\arabic{enumi}.}
\setcounter{enumi}{1}
\tightlist
\item
  Reformulación como producto cartesiano: hacemos el producto cartesiano
  entre los buckets. Como solo hay una posibilidad por bucket solo
  podremos hacer una reformulación. Si nos fijamos, esta reformulación
  será la misma que la reformulación GAV:
\end{enumerate}

\begin{verbatim}
Q1''(id, nombre_estacion, nombre_municipio):- data_renfe_ESTACION(v1, id, nombre_estacion, 
lat_e, lon_e, v2, v3, nom_mun, v4, v5), wikidata_MUNICIPIO(id_mun, nombre_municipio, v6, v7), 
wikidata_MUNICIPIO(id_mun, nombre_municipio, hab, coord), lat_m=latitud(coord), 
lon_m=longitud(coord), INE_MUNICIPIO(cmun, cpro, v8, v9), 
INE_PROVINCIA(cpro, prov, v10, ccaa), id_mun=cmun+cpro, hab<10000
\end{verbatim}

\begin{enumerate}
\def\labelenumi{\arabic{enumi}.}
\setcounter{enumi}{2}
\tightlist
\item
  Posible simplificación: como ya vimos anteriormente, es posible
  simplificar la consulta, por lo que el resultado final es:
\end{enumerate}

\begin{verbatim}
Q1'(id, nombre_estacion, nombre_municipio):- data_renfe_ESTACION(v1, id, nombre_estacion, 
lat_e, lon_e, v2, v3, nom_mun, v4, v5), wikidata_MUNICIPIO(id_mun, nombre_municipio, hab, 
coord), lat_m=latitud(coord), lon_m=longitud(coord), INE_MUNICIPIO(cmun, cpro, v8, v9),
INE_PROVINCIA(cpro, prov, v10, ccaa), id_mun=cmun+cpro, hab<10000
\end{verbatim}

Como ya hemos visto, la reformulación LAV y GAV coinciden aunque hayamos
utilizado diferentes algoritmos para su construcción. Es por esto que
para las siguientes dos consultas únicamente se dejará indicada la
reformulación final.

\subsection{Consulta 2}\label{consulta-2}

Esta consulta tiene un \texttt{OR} para poder trabajar con la tabla
distancia. Es por ello que la consulta se define como:

\[
Q2=Q2_a\cup Q2_b
\]

Estas subconsultas se pueden escribir sobre el esquema mediador como:

\begin{verbatim}
Q2a(id):- Viaje(id, ruta, v1, v2), Ruta(ruta, origen, destino, v3), 
Distancia(origen, destino, dist), dist<30
\end{verbatim}

\begin{verbatim}
Q2a(id):- Viaje(id, ruta, v1, v2), Ruta(ruta, origen, destino, v3), 
Distancia(destino, origen, dist), dist<30
\end{verbatim}

\subsubsection{Reformulación GAV/LAV}\label{reformulaciuxf3n-gavlav}

La consulta en forma GAV queda extremadamente larga, puesto que para
calcular distancias hay que cruzar datos con la vista auxiliar

\begin{verbatim}
Q2a'(id):- Parada_aux(ruta, num_secuencia, nombre_parada), 
nombre_origen=getMinSecuencia(num_secuencia), data_renfe_ESTACION(v1, estacion, nombre_origen, 
v3, v4, v5, v6, v7, v8, pais), adif_SALIDAS(fecha, horario, v10, v11, estacion), 
id=secuencia(), pais==España, Parada_aux(ruta, num_secuencia, nombre_parada), 
nombre_origen=getMinSecuencia(num_secuencia), nombre_destino=getMaxSecuencia(num_secuencia), 
data_renfe_ESTACION(v1, origen, nombre_origen, v12, v13, v14, v15, v16, v17, pais), 
data_renfe_ESTACION(v36, destino, nombre_destino, v18, v19, v20, v21, v22, v23, pais), 
pais==España, data_renfe_ESTACION(v24, origen, v25, latitud_1, longitud_1, v26, v27, 
v28, v29, pais), data_renfe_ESTACION(v30, destino, v31, latitud_2, longitud_2, v32, 
v33, v34, v35, pais), distancia=dist(latitud_1, latitud_2, longitud_1, longitud_2), 
origen < destino, pais==España, dist<30
\end{verbatim}

Tras simplificar, la consulta queda como:

\begin{verbatim}
Q2a'(id):- Parada_aux(ruta, num_secuencia, nombre_parada), 
nombre_origen=getMinSecuencia(num_secuencia), nombre_destino=getMaxSecuencia(num_secuencia), 
data_renfe_ESTACION(v1, origen, nombre_origen, latitud_1, longitud_1, v2, v3, v4, v5, v6), 
data_renfe_ESTACION(v7, destino, nombre_destino, latitud_2, longitud_2, v8, v9, v10, v11, 
v12), adif_SALIDAS(fecha, horario, v13, v14, origen), id=secuencia(), 
distancia=dist(latitud_1, latitud_2, longitud_1, longitud_2), 
pais==España, origen < destino, distancia<30
\end{verbatim}

Q2b' cubrirá aquellos casos en los que destino \textless{} origen.

\begin{verbatim}
Q2a'(id):- Parada_aux(ruta, num_secuencia, nombre_parada), 
nombre_origen=getMinSecuencia(num_secuencia), nombre_destino=getMaxSecuencia(num_secuencia), 
data_renfe_ESTACION(v1, origen, nombre_origen, latitud_1, longitud_1, v2, v3, v4, v5, v6), 
data_renfe_ESTACION(v7, destino, nombre_destino, latitud_2, longitud_2, v8, v9, v10, v11, 
v12), adif_SALIDAS(fecha, horario, v13, v14, origen), id=secuencia(), 
distancia=dist(latitud_1, latitud_2, longitud_1, longitud_2), 
pais==España, destino < origen, distancia<30
\end{verbatim}

\subsection{Consulta 3}\label{consulta-3}

Si vemos la definición de la consulta 3, podemos ver que para poder
obtener las estaciones con 0 viajes entre sí debemos hacer un
\texttt{LEFT\ JOIN} sobre el esquema mediador, y no sabemos como poner
dicha operación en forma conjuntiva. Por tanto, tampoco sabríamos
escribir la consulta con la reformulación GAV o la reformulación LAV.

\section{Scraping}\label{scraping}

\subsection{ADIF}\label{adif}

Desde la página web de consulta de horarios de estaciones en tiempo real
de\href{https://www.adif.es/viajeros/estaciones}{ADIF} hemos extraido
mediante scrapping las tablas de las fuentes adif\_SALIDAS y
adif\_LLEGADAS. No ha sido sencillo, pues es una aplicación web
relativamente moderna y está bien preparada frente a interacción
atomatizada, lo cuál ha resultado un reto. La implementación del
scraping se puede ver en el archivo `scraping\_adif.py'. Este fichero
actua a modo demo, por lo que no hay ningún tipo de interacción con el
usuario. Simplemente accede a una página (la estación de Valladolid) y
extrae las dos tablas. A continuación se procede a explicar la lógica y
detalles técnicos:

Una petición http básica no era suficiente para poder acceder al
contenido real, pues la aplicación respondía con un html especial
indicando que el acceso no estaba permitido:

\begin{Shaded}
\begin{Highlighting}[]
\DataTypeTok{\textless{}}\KeywordTok{HTML}\DataTypeTok{\textgreater{}}
\DataTypeTok{\textless{}}\KeywordTok{HEAD}\DataTypeTok{\textgreater{}}
\DataTypeTok{\textless{}}\KeywordTok{TITLE}\DataTypeTok{\textgreater{}}\NormalTok{Access Denied}\DataTypeTok{\textless{}/}\KeywordTok{TITLE}\DataTypeTok{\textgreater{}}
\DataTypeTok{\textless{}/}\KeywordTok{HEAD}\DataTypeTok{\textgreater{}}
\DataTypeTok{\textless{}}\KeywordTok{BODY}\DataTypeTok{\textgreater{}}
\DataTypeTok{\textless{}}\KeywordTok{H1}\DataTypeTok{\textgreater{}}\NormalTok{Access Denied}\DataTypeTok{\textless{}/}\KeywordTok{H1}\DataTypeTok{\textgreater{}}

\NormalTok{You don\textquotesingle{}t have permission to access}
\DataTypeTok{\textless{}/}\KeywordTok{HTML}\DataTypeTok{\textgreater{}}
\end{Highlighting}
\end{Shaded}

por lo que tuvimos que rechazar una primera aproximación usando los
módulos \texttt{requests} con \texttt{BeautifullSoup}. Por tanto,
probamos a usar el módulo \texttt{playwright}, desarrollado por
Microsoft y mucho más potente.

El módulo proporciona una interfáz muy sencilla para manejar versiones
``headless'' de un navegador, con el objetivo de simular la navegación
de un usuario real, pudiendo interactuar con elementos de la aplicación
sin simplemente limitarse a la respuesta de una petición básica. Aún
así, el sistema de detección de ``bots'' de la web seguía respondiendo
con ficheros html de rechazo, pues parece ser que también era capaz de
identificar que se trataba de un navegador ``headless''. Pudimos
solucionar este problema creando un nuevo ``contexto'' de navegador que
simulase una ventana gráfica, pero sin que de verdad esta existiese;
gracias a esto hemos sido capaces de acceder a la información.

El flujo del programa es muy simple:

\begin{enumerate}
\def\labelenumi{\arabic{enumi}.}
\tightlist
\item
  Se crea una ``ventana'' virtual y una página sobre esta,
  posteriormente se navega a la ruta (para el caso de valladolid,
  \url{https://www.adif.es/w/10600-valladolid-c.-g.}) y se espera a que
  todo el DOM esté en estado \emph{attached} para asegurar que los
  elementos a los que queremos acceder están disponibles en el momento
  de buscarlos.
\item
  Dependiendo de si se quieren estraer las tablas de llegadas o salidas,
  hay que pulsar un botón de ``salidas'' o no (la opción de ``llegadas''
  parece que viene seleccionada por defecto). Este es el único paso
  diferente, a partir de aquí el programa es igual para ambos salvo
  alguna pequeña variante de notación. Para acceder a la tabla hay que
  filtrar por el identificador
  \texttt{\#horas-trenes-estacion-(salidas\textbar{}llegadas)}
  dependiendo de cual sea la fuente que se quiera extraer.
\end{enumerate}

\begin{figure}
\centering
\includegraphics{./img/img1_scrapping_ADIF.png}
\caption{Captura de pantalla de la página de ADIF}
\end{figure}

\begin{enumerate}
\def\labelenumi{\arabic{enumi}.}
\setcounter{enumi}{2}
\tightlist
\item
  La extracción de información es la misma, pues usa los mismos
  identificadores para las columnas, incluso la estación de origen o
  destino también tiene el mismo identificador \texttt{col-destino}
  independientemente de si se habla de salidas o llegadas. Al ser una
  tabla, extraer los datos es tan sencillo como iterar sobre todos los
  objetos de tipo \texttt{tr} y acceder a cada elemento \texttt{td} por
  columna. Una cosa importante y que tuvimos que solucionar fue que hay
  dos tipos de ``filas'', las que contienen la información que nos
  interesa y las amarillas, que proporcionan información sobre problemas
  en la infrastructura, que no tienen el mismo formato y por tanto
  pueden producir errores si no se trata bien.
\item
  Finalmente, se ``serializan'' los datos extraidos en un csv usando la
  función de la librería estandar \texttt{csv.DictWritter}.
\end{enumerate}

\subsection{RENFE}\label{renfe}

Hemos utilizado scraping sobre la web de
\href{https://www.renfe.com/es/es/viajar/informacion-util/horarios}{renfe
horarios}. Para conseguir esta información hemos tenido que rellenar un
formulario utilizando \emph{playwright} y \emph{BeautifulSoup} en Python
para acceder, ya que con la librería \emph{requests} aparecían problemas
para realizar el scraping.

El scraping realizado se puede ver en el archivo
\emph{scraping\_renfe\_horarios.py}, y el flujo del programa es el
siguiente:

\begin{enumerate}
\def\labelenumi{\arabic{enumi}.}
\tightlist
\item
  Se introduce al programa la información relativa al cuestionario con
  el siguiente formato:
  {[}cod\_estacion\_origen{]}-{[}cod\_estacion\_destino{]}-{[}dia{]}-{[}mes{]}-{[}año{]},
  de tal forma que para obtener los recorridos entre Valladolid Campo
  Grande (estación número 10600) y Palencia (estación número 14100) para
  el día 20 de abril de 2026 podríamos ejecutar el programa como:
\end{enumerate}

\texttt{\{bash\}\ \$\ echo\ "10600-14100-20-04-2026"\ \textbar{}\ python3\ scraping\_renfe\_horarios.py}

Como se puede ver, el programa solo hace scraping sobre dos estaciones
para un día concreto. El motivo de esto es que todavía no heos obtenido
el listado de todas las estaciones para poder hacer el scraping de todos
los viajes, aunque el programa está preparado para soportar varias
llamadas seguidas.

\begin{enumerate}
\def\labelenumi{\arabic{enumi}.}
\setcounter{enumi}{1}
\item
  Se accede a la web, accediendo al \emph{iframe} que contiene el
  formulario. Si no se accede a este iframe no se pueden obtener los
  elementos html del formulario. Cuando se accede, se rellenan los datos
  acorde a la entrada del programa y se hace click en el botón
  \textbf{BUSCAR}.
\item
  El \emph{iframe} cambia y muestra una tabla como la de la imagen de
  debajo. Accedemos a cada linea de la tabla y obtenemos el primer
  valor, que redirige a otra página web mediante una función javaScript
  y un hiperenlace.
\end{enumerate}

\begin{figure}
\centering
\includegraphics{./img/img1_scraping_RENFE.png}
\caption{Tabla de horarios RENFE}
\end{figure}

\begin{enumerate}
\def\labelenumi{\arabic{enumi}.}
\setcounter{enumi}{3}
\tightlist
\item
  Formateamos el hiperenlace para hacerlo coincidir con el correcto.
  Esto equivale a eliminar saltos de línea y sustituir espacios en
  blanco por \emph{\%20}, de tal forma que la url:
\end{enumerate}

\begin{Shaded}
\begin{Highlighting}[]
\NormalTok{horarios.renfe.com/HIRRenfeWeb/recorrido.do?O=10600}\ErrorTok{\&}\NormalTok{D=14100}\ErrorTok{\&}\NormalTok{F=2026{-}04{-}20}\ErrorTok{\&}\NormalTok{T=04073}\ErrorTok{\&}\NormalTok{G=1}\ErrorTok{\&}\NormalTok{TT=ALVIA\%20\%20\%20\%20\%20\%20\%20\%20\%20\%20\%20\%20\%20\%20 }\ErrorTok{\&}\NormalTok{ID=s}\ErrorTok{\&}\NormalTok{FDS=2026{-}04{-}20}\ErrorTok{\&}\NormalTok{DT=33 min.}
\end{Highlighting}
\end{Shaded}

Equivale a la imagen de debajo.
\includegraphics{./img/img2_scraping_RENFE.png}

\begin{enumerate}
\def\labelenumi{\arabic{enumi}.}
\setcounter{enumi}{4}
\item
  Accedemos a estas url (una por ruta) y obtenemos el listado de paradas
\item
  Accedemos a dos archivos csv, correspondientes a los Esquemas en
  Origen de \emph{horarios\_renfe\_RUTA} y
  \emph{horarios\_renfe\_HORARIO}, y rellenamos los datos. Accedemos a
  estos archivos en modo append: si no existen se crean y si existen se
  añade la nueva información al final (así podemos crear muchas llamadas
  al programa para diferentes rutas).
\end{enumerate}

Por último, se ha añadido la opción \emph{--verbose} para obtener logs
del programa.

\section{SPARQL}\label{sparql}

\subsection{WikiData}\label{wikidata}

Los datos de los municipios de España los vamos a coger de
\includegraphics{https://www.wikidata.org/?uselang=es}, mediante un
script de \emph{Python} que utilice el servicio de consultas del propio
WikiData.

En este caso, la consulta ha sido directa y sencilla, necesitando
únicamente el nombre del municipio, la población que tiene y sus
coordenadas, en formato POINT. Únicamente hay que tener en cuenta que
WikiData puede cancelar la consulta si no se especifica un User-Agent
(debido a su política de Robots).

\begin{longtable}[]{@{}ll@{}}
\toprule\noalign{}
Columna & Tipo \\
\midrule\noalign{}
\endhead
\bottomrule\noalign{}
\endlastfoot
CODIGO & text \\
LABEL & text \\
POOBLACION & numeric \\
COORDENADAS & Point \\
\end{longtable}

Para lograrlo, hemos hecho la siguiente consulta:

\begin{Shaded}
\begin{Highlighting}[]
\KeywordTok{SELECT}\NormalTok{ ?codigo ?}\KeywordTok{label}\NormalTok{ ?poblacion ?coordenadas }\KeywordTok{WHERE}\NormalTok{ \{}
\NormalTok{  ?municipio (wdt}\CharTok{:P31}\OperatorTok{/}\NormalTok{(wdt}\CharTok{:P279}\OperatorTok{*}\NormalTok{)) wd}\CharTok{:Q2074737}\NormalTok{;}
\NormalTok{    wdt}\CharTok{:P772}\NormalTok{ ?codigo;}
\NormalTok{    wdt}\CharTok{:P1082}\NormalTok{ ?poblacion;}
\NormalTok{    wdt}\CharTok{:P625}\NormalTok{ ?coordenadas.}

\NormalTok{  SERVICE wikibase}\CharTok{:label}\NormalTok{ \{}
\NormalTok{    bd}\CharTok{:serviceParam}\NormalTok{ wikibase}\CharTok{:language} \OtherTok{"es"}\NormalTok{.}
\NormalTok{    ?municipio rdfs}\CharTok{:label}\NormalTok{ ?}\KeywordTok{label}\NormalTok{.}
\NormalTok{  \}}
\NormalTok{\}}
\KeywordTok{ORDER} \KeywordTok{BY}\NormalTok{ ?codigo}
\end{Highlighting}
\end{Shaded}
